\documentclass[journal,twoside,web]{ieeecolor}

\usepackage{generic}
\usepackage{cite}
\usepackage{amsmath,amssymb,amsfonts}
\usepackage{graphicx}
\usepackage{multirow}
\usepackage{booktabs}
\usepackage{algorithm,algorithmic}
\usepackage{hyperref}
\hypersetup{hidelinks}
\usepackage{textcomp}
\usepackage{placeins}
\def\BibTeX{{\rm B\kern-.05em{\sc i\kern-.025em b}\kern-.08em
    T\kern-.1667em\lower.7ex\hbox{E}\kern-.125emX}}

\begin{document}

\title{\vspace{15mm}WavKAN-v2: Quantified Structural Interpretability and Curriculum-Anchored Minority-Class Recall for Inter-Patient ECG Arrhythmia Classification}

\author{Venkateramanan~Manivannan and Ramanathan~Lakshmanan
\thanks{Manuscript submitted for review to the IEEE Journal of Biomedical and Health Informatics.}
\thanks{Venkateramanan Manivannan is with the School of Computer Science
and Engineering (SCOPE), Vellore Institute of Technology (VIT), Vellore,
Tamil Nadu 632014, India (e-mail: venkate.ramanan2024@vitstudent.ac.in).}
\thanks{Ramanathan Lakshmanan is a Professor with the Department of IoT,
School of Computer Science and Engineering, Vellore Institute of Technology (VIT), Vellore,
Tamil Nadu 632014, India (Corresponding author, e-mail: lramanathan@vit.ac.in).}}

\markboth{\hskip25pc IEEE JOURNAL OF BIOMEDICAL AND HEALTH INFORMATICS}
{Venkateramanan Manivannan \MakeLowercase{\textit{et al.}}: WavKAN-v2 for Inter-Patient Arrhythmia Detection}

\maketitle

\begin{abstract}
Automated arrhythmia classification from ECG signals is limited by two persistent, under-addressed problems: the opacity of high-performing black-box classifiers, and the tendency of reported metrics to be measured under evaluation protocols that leak patient identity between training and test sets. We present WavKAN-v2, a Wavelet Kolmogorov--Arnold Network (154,325 parameters) that replaces fixed activation functions with learnable Mexican Hat wavelet basis functions, combined with a Progressive Curriculum Anchoring (PCA) training schedule for severe class imbalance, and evaluated exclusively under the strict inter-patient DS1/DS2 protocol (de Chazal et al.) on the MIT-BIH Arrhythmia Database. Under a rigorous 20-seed evaluation with Holm-Bonferroni-corrected Wilcoxon signed-rank tests and Cohen's $d$ effect sizes, we report two honest, load-bearing findings rather than a single accuracy claim. First, curriculum learning produces a real, statistically robust improvement in Supraventricular (S) recall (0.079$\pm$0.016 $\to$ 0.123$\pm$0.052, Holm-adjusted $p=0.0016$, $d=-0.77$) with no statistically significant cost to Macro-F1, V-recall, N-recall, or F-recall. Second, and reported without narrative softening, WavKAN-v2's Macro-F1 (0.323$\pm$0.018) is significantly \emph{lower} than four standard, substantially smaller architectures evaluated under an identical protocol and training budget --- ResNet1D (62,869 params, 0.351$\pm$0.018), a lightweight Transformer (29,653 params, 0.353$\pm$0.039), a CNN with Focal Loss (71,061 params, 0.362$\pm$0.015), and a B-Spline KAN of the same architectural family (118,229 params, 0.371$\pm$0.023) --- all four comparisons surviving Holm correction. Ventricular (V) recall, by contrast, is comparable across all five models (0.88--0.91), indicating the gap is Macro-F1-specific rather than a uniform failure. What WavKAN-v2 offers instead of a raw-accuracy advantage is a quantified interpretability metric unavailable to any comparator: a Wavelet--ECG Alignment Score (WAS) of 0.971, obtained by directly measuring the correspondence between each learned wavelet's translation/dilation parameters and known ECG fiducial priors. We additionally report real, measured CPU inference latency (3.18$\pm$1.22 ms/beat, batch size 1) and an honest negative result for post-training INT8 quantization, which provides no benefit for this architecture (1.08$\times$ size reduction, 28\% \emph{slower} inference) because dynamic quantization only touches \texttt{nn.Linear} layers, and this model's parameter mass sits in the custom KAN layer and a bidirectional GRU. Zero-shot cross-dataset transfer to PTB-XL (single seed, 75,170 beats) remains weak (Macro-F1 0.204), which we report as an open limitation rather than omit. We position WavKAN-v2's contribution accordingly: not as an accuracy-superior replacement for standard architectures, but as an intrinsically, quantifiably interpretable alternative with a real, targeted minority-class recall benefit, evaluated with a level of statistical rigor (20 seeds, multiple-comparison correction, effect sizes) that is uncommon in the inter-patient ECG classification literature.
\end{abstract}

\begin{IEEEkeywords}
Kolmogorov-Arnold Networks, Wavelet, Arrhythmia Classification, Quantified Interpretability, Curriculum Learning, Inter-Patient Generalization, Statistical Rigor
\end{IEEEkeywords}

\section{Introduction}

\IEEEPARstart{C}{ardiovascular} diseases (CVDs) are currently the leading cause of death worldwide and require reliable and continuous cardiac monitoring. While the Electrocardiogram (ECG) is regarded as the gold standard for the detection of arrhythmias, manual interpretation of long-term Holter recordings is labor intensive, error-prone and infeasible for large-scale screening. To automate this process, deep learning architectures, in particular Convolutional Neural Networks (CNNs) \cite{takalo-mattila_inter-patient_2018, guo_inter-patient_2019} and Transformers \cite{akan_ecgformer_2024} have been widely adopted with high accuracy under controlled experimental conditions.

However, despite these improvements in accuracy, existing state-of-the-art solutions encounter three major hurdles to safe clinical deployment, and, as we show in this paper, a fourth methodological hurdle that is rarely confronted directly: the difficulty of reporting an honest, statistically rigorous comparison once a genuinely fair inter-patient protocol is enforced on all sides of the comparison.

\begin{enumerate}
    \item \textbf{Opacity of Black-Box Models:} The lack of inherent interpretability in CNNs and Transformers limits the ability of clinicians to associate learned representations with underlying physiological markers, such as P-wave morphology. This lack of interpretability reduces clinical trust and complicates regulatory approval \cite{taleban_explainable_2026}.
    \item \textbf{Inter-Patient Generalization Gap:} A major limitation of previous studies is intra-patient evaluation, in which beats from the same patient are mixed across training and test sets, inflating reported performance via patient-specific overfitting \cite{bahrami_investigation_2025}. Under strict inter-patient evaluation, models are often shown to have significant performance degradation, particularly for minority arrhythmia classes \cite{xiao_deep_2023, silva_systematic_2025}.
    \item \textbf{Computational Inefficiency:} Architectures such as Vision Transformers (ViT) require millions of parameters, making them unsuitable for battery-powered wearable devices \cite{elsheikhy_lightweight_2025}.
    \item \textbf{Unverified Interpretability and Uncorrected Statistics:} ``Interpretable'' is frequently asserted qualitatively (a figure of learned filters) rather than measured quantitatively against a physiological ground truth, and multi-metric comparisons are frequently reported without correcting for multiple comparisons, inflating the apparent rate of significant findings.
\end{enumerate}

To address these four issues jointly, i.e.\ interpretability (both qualitative and quantitative), inter-patient generalization, edge efficiency, and statistical honesty, we explore the use of Kolmogorov-Arnold Networks (KANs) \cite{liu_kan_2024} as a structural foundation. Unlike standard Multi-Layer Perceptrons, KANs use learnable activation functions along edges instead of fixed nodes, making the architecture structurally transparent. However, most existing KAN-based ECG models, such as MAK-Net \cite{zhao_mak-net_2025}, rely on B-spline basis functions, which may not optimally capture sharp transient components such as the QRS complex, and none of the KAN-ECG literature we are aware of quantifies the resulting interpretability claim against a physiological reference.

To better capture sharp transient ECG morphology, we substitute B-splines with learnable Mexican Hat Wavelets \cite{addison_wavelet_2005, bozorgasl_liu_2024}. To address class imbalance, we use a two-phase, minority-anchored curriculum (Progressive Curriculum Anchoring, PCA) \cite{bengio_curriculum_2009, schmale_curriculum_2025}. Critically, and unlike prior applications of curriculum learning to this problem, we evaluate its effect with 20-seed, Holm-Bonferroni-corrected statistics against a hyperparameter-matched fair baseline, so that the reported effect is not an artifact of uncorrected multiple comparisons or a mismatched training recipe.

The contributions of this study are:
\begin{enumerate}
    \item \textbf{A quantified interpretability metric.} We introduce the Wavelet--ECG Alignment Score (WAS), which measures the numerical correspondence between each learned wavelet edge's translation ($\mu$) and dilation ($\gamma$) parameters and known QRS/P/T-wave priors. WavKAN-v2 achieves a Macro-WAS of 0.971 --- to our knowledge the first quantitative (not solely visual) interpretability metric reported for a KAN-based ECG classifier.
    \item \textbf{A statistically rigorous, honestly reported cross-architecture comparison.} Under an identical protocol, data split, and training budget, we compare WavKAN-v2 against four standard architectures (ResNet1D, a lightweight Transformer, a CNN with Focal Loss, and a same-family B-Spline KAN) across 20 seeds each, with Holm-Bonferroni-corrected Wilcoxon signed-rank tests and Cohen's $d$. We report the result plainly: all four baselines significantly outperform WavKAN-v2 on Macro-F1, despite WavKAN-v2 being the largest model of the five.
    \item \textbf{A real, targeted curriculum-learning benefit.} Progressive Curriculum Anchoring produces a statistically robust improvement in Supraventricular (S) recall with no significant cost to any other metric, once measured with correction for multiple comparisons across the five metrics tested.
    \item \textbf{Honest efficiency and generalization reporting.} We report real, measured CPU latency and an honest negative quantization result, and we report weak (not strong) zero-shot cross-dataset transfer to PTB-XL rather than omitting it.
\end{enumerate}

%% 2. RELATED WORK
\section{Related Work}
\label{sec:related_work}

The automated classification of ECG arrhythmias has changed significantly in the past decade. However, the literature shows consistent methodological differences in evaluation protocol, model transparency, and computational efficiency, and --- as this paper's own results illustrate --- an under-examined tendency for multi-metric comparisons to go unreported without correction for multiple comparisons.

\subsection{Inter-Patient Evaluation Paradigm}
One of the key methodological differences in ECG classification is the choice between intra-patient and inter-patient evaluation. As underscored by systematic reviews by Xiao et al. \cite{xiao_deep_2023} and Silva et al. \cite{silva_systematic_2025}, many early deep learning studies adopted the intra-patient paradigm in which heartbeats from a single patient are randomly allocated for training and testing sets. While this approach results in extremely high accuracies ($> 97\%$), it leads to leakage of data, as models tend to memorize the morphological baselines for each individual patient, rather than trying to generalize to physiological markers of arrhythmia \cite{bahrami_investigation_2025}. To address this, De Chazal et al. \cite{chazal_automatic_2004} proposed the strict inter-patient protocol (DS1/DS2 split) for the MIT-BIH dataset, where patients in the test set are entirely unseen. Recent state-of-the-art architectures evaluated under this protocol include the convolutional and recurrent neural network (CNN+RNN) by Guo et al. \cite{guo_inter-patient_2019}, the multiscale CNN with attention by Zhou et al. \cite{zhou_inter-patient_2024}, and the sequence-to-sequence models by Mousavi et al. \cite{mousavi_inter-_2019}.

\subsection{Explainability and Kolmogorov-Arnold Networks}
The black-box nature of standard CNN and Transformer architectures is a major challenge to clinical adoption. A recent systematic review by Taleban et al. \cite{taleban_explainable_2026} underscored the need to move from post-hoc explanation methods (e.g., Grad-CAM, LIME), which can be computationally expensive and sensitive to noise, towards ``intrinsically interpretable'' models.

Kolmogorov-Arnold Networks (KANs) \cite{liu_kan_2024} place learnable univariate activation functions on network edges rather than fixed node activations. Early applications of KANs to ECG signals, such as MAK-Net by Zhao et al. \cite{zhao_mak-net_2025}, showed promising accuracy using B-spline bases, but B-splines may have difficulty separating fast transient physiological features \cite{bozorgasl_liu_2024}. Bozorgasl and Chen \cite{bozorgasl_liu_2024} proposed Wav-KAN, using wavelet functions instead, consistent with classical ECG time-frequency analysis theory. Per Addison's theory of Mexican Hat wavelet analysis \cite{addison_wavelet_2005}, the Mexican Hat wavelet has a natural representation of QRS time-frequency geometry, making it a physiologically motivated basis for interpretable feature extraction. However, to our knowledge no prior KAN-ECG study has \emph{quantified} the resulting interpretability against a physiological ground truth, or evaluated it with corrected multi-seed statistics against a hyperparameter-matched set of standard baselines --- both of which this paper does directly.

\subsection{Handling Imbalance and Edge Constraints}
Inter-patient MIT-BIH datasets show extreme imbalance ratios (Normal to Fusion beats exceeding 100:1) \cite{silva_systematic_2025}. Standard cross-entropy training often leads to majority-class collapse. Curriculum Learning \cite{bengio_curriculum_2009} takes a dynamic approach in which the sampling distribution or loss weighting changes during training. Bae \& Shin \cite{bae_handling_2025} showed that naive contrastive learning on unbalanced medical data can \emph{degrade} performance, while a curriculum-phased, distribution-aware variant restored and improved it. Schmale et al. \cite{schmale_curriculum_2025} similarly showed that curricula prioritizing rare, complex phenotypes can stabilize decision boundaries. These findings motivate our own minority-anchored curriculum --- but, as our results show (\S\ref{sec:results}), motivation is not the same as an accuracy guarantee: we find the curriculum's real, statistically defensible benefit is narrower (S-recall specifically) than the broad ``improves minority-class detection'' framing common in this literature.

Deploying these models in wearable Holter monitors also requires attention to computational cost. Farag \cite{farag_tiny_2023} and Elsheikhy et al. \cite{elsheikhy_lightweight_2025} emphasize reduced memory footprints for on-device inference; we report real measured latency and quantization behavior for WavKAN-v2 specifically (\S\ref{sec:results}), rather than an unmeasured efficiency claim.

%% 3. METHODOLOGY
\section{Methodology}
\label{sec:methods}

\subsection{Problem Formalization}
The ECG signal is segmented into beat-centered windows $x_i \in \mathbb{R}^T$ ($T=360$ samples at 360\,Hz) together with a causal RR-interval history $r_i \in \mathbb{R}^K$ ($K=5$, comprising only the current and four \emph{preceding} inter-beat intervals, never future intervals, so that the representation is compatible with real-time inference) to learn $f:(x_i,r_i)\to y_i$, $y_i \in \{N,S,V,F,Q\}$, under severe class imbalance ($N \gg V > S > F \gg Q$).

\subsection{The WavKAN-v2 Architecture}
The proposed system (Fig.~\ref{fig:arch}) employs a dual-branch architecture combining a learnable wavelet backbone (morphology) with a rhythm-aware branch (RR history), fused before classification.

\begin{figure*}[!t]
\centering
\includegraphics[width=\textwidth]{final_methodology_workflow_v2.pdf}
\caption{\textbf{WavKAN-v2 Architecture.} Dual-branch architecture combining the morphological WavKAN backbone with an RR-history rhythm branch. The curriculum scheduler (training only) anchors minority-class decision boundaries before annealing to the natural class distribution.}
\label{fig:arch}
\end{figure*}

\subsubsection{Feature Extraction: WavKAN Backbone}
Morphological features $z_m$ are extracted by a dense Wavelet Kolmogorov-Arnold Network (WavKAN) layer with physiology-constrained wavelet initialization (PCWI) and a P-wave attention module (PWAM). Each edge learns an independently parameterized Mexican Hat wavelet $\psi(t;\mu,\gamma)$, the negative normalized second derivative of a Gaussian, chosen because it behaves as a high-frequency transient detector well matched to the sharp QRS complex, unlike smooth B-spline polynomials or the oscillatory Morlet wavelet \cite{addison_wavelet_2005}. The WavKAN layer maps $x_i \in \mathbb{R}^{360}$ to $z_{kan} \in \mathbb{R}^{64}$:
\begin{equation}
z_{kan}^{(k)} = \sum_{j=1}^{360} w_{j,k} \cdot \psi\left(\frac{x_j - \mu_{j,k}}{\gamma_{j,k}}\right), \quad k=1,\dots,64.
\end{equation}
The output is normalized (LayerNorm) and passed through a Bidirectional GRU (32 hidden units/direction) to produce $z_m \in \mathbb{R}^{64}$.

\begin{figure*}[!t]
\centering
\includegraphics[width=0.9\textwidth]{wavkan_micro_architecture_v2.pdf}
\caption{\textbf{Micro-Architecture of WavKAN.} Learnable wavelet functions $\psi(t;\mu,\gamma)$ replace fixed node activations, with each edge learning its own translation and dilation parameters.}
\label{fig:wavkan_micro}
\end{figure*}

\subsubsection{Rhythm Encoding}
The causal RR-history vector $R_{seq} = [RR_t, RR_{t-1}, RR_{t-2}, RR_{t-3}, RR_{t-4}]$ is normalized by the local moving average of the preceding 10 beats, then passed through a self-attention branch and a small MLP to produce $z_r \in \mathbb{R}^{16}$. \S\ref{sec:results} reports a within-family ablation (Table~\ref{tab:family_ablation}) showing this RR self-attention branch is, on the current data, \emph{not} the best-performing choice among the variants we tested --- a finding we report and discuss (\S\ref{sec:discussion}) rather than omit.

\subsubsection{Fusion and Classification}
$z=[z_m;z_r]\in\mathbb{R}^{80}$ is passed to a final classifier ($80\to48\to5$). Total trainable parameters: \textbf{154,325}, a $>$95\% reduction relative to MAK-Net (6.1M) \cite{zhao_mak-net_2025} and DenseNet+GRU architectures \cite{guo_inter-patient_2019}.

\subsection{Progressive Curriculum Anchoring (PCA)}
Unlike a generic easy-to-hard curriculum, PCA restructures training by class rarity, in two phases:
\begin{enumerate}
    \item \textbf{Warm-up} ($\approx$25\% of the epoch budget): a fully class-balanced \texttt{WeightedRandomSampler} (every class equally likely per batch) with plain, unweighted cross-entropy, so that no class's decision boundary is anchored purely by frequency before the others are seen.
    \item \textbf{Anneal} (remaining epochs): the sampler's balancing exponent decays linearly from 1.0 (fully balanced) to 0.0 (natural distribution), paired with Focal Loss ($\gamma=2.0$, inverse-frequency class weights, with the S-class weight further multiplied by $8\times$ to counteract its morphological similarity to N).
\end{enumerate}
This schedule is training-only; inference uses the natural class distribution. The \textbf{fair-baseline arm} used for every comparison in \S\ref{sec:results} is architecturally identical and uses natural-distribution sampling for every epoch with plain class-weighted cross-entropy (same weight formula, no warm-up/anneal distinction) --- i.e., the only variable that differs between arms is the sampling/loss schedule itself, not the architecture, optimizer, or epoch budget.

\begin{algorithm}[!t]
\caption{Progressive Curriculum Anchoring (PCA)}
\label{alg:curriculum}
\begin{algorithmic}[1]
\STATE \textbf{Input:} Training set $\mathcal{D}$, total epochs $E$, warm-up fraction $\alpha=0.25$
\STATE \textbf{Output:} Trained model $f_\theta$
\STATE Compute class frequencies $\{n_c\}$, $c \in \{N,S,V,F,Q\}$; inverse-frequency weights $w_c$; $w_S \gets 8\cdot w_S$
\FOR{epoch $e=1$ to $E$}
\IF{$e \le \alpha E$}
\STATE Sample mini-batches via a fully class-balanced sampler (balance exponent $\beta=1.0$)
\STATE $\mathcal{L} \gets \mathcal{L}_{\text{CE}}$ (unweighted)
\ELSE
\STATE $\beta \gets$ linear decay from 1.0 to 0.0 over the remaining epochs
\STATE Sample mini-batches via sampler with balance exponent $\beta$
\STATE $\mathcal{L} \gets \text{FocalLoss}(\gamma=2.0, w_c)$
\ENDIF
\STATE Update $\theta$ via AdamW on $\mathcal{L}$; checkpoint on validation Macro-F1 improvement
\ENDFOR
\end{algorithmic}
\end{algorithm}

%% DATASET
\section{Dataset and Evaluation Protocol}

\subsection{Inter-Patient Data Splitting Protocol}
We follow the inter-patient paradigm of De Chazal et al. \cite{chazal_automatic_2004}. The MIT-BIH Arrhythmia Database is split into DS1 (train+validation) and DS2 (test), as in Table~\ref{tab:split_protocol}; DS1 is further split into training and validation subsets, following Takalo-Mattila et al. \cite{takalo-mattila_inter-patient_2018}, so that early stopping and model selection never touch DS2. All test results in \S\ref{sec:results} are computed exclusively on DS2.

\begin{table*}[!t]
\caption{Inter-Patient Data Splitting Protocol (DS1/DS2).}
\label{tab:split_protocol}
\centering
\footnotesize
\begin{tabular}{@{}l l p{6.5cm} l@{}}
\toprule
\textbf{Set} & \textbf{Role} & \textbf{Record IDs (MIT-BIH)} & \textbf{Function} \\ \midrule
\multirow{2}{*}{\textbf{DS1}} & \textbf{Train} (18 rec.) & 101, 106, 108, 109, 112, 114, 115, 116, 118, 119, 122, 124, 201, 203, 205, 207, 215, 220 & Optimization \\
& \textbf{Val} (4 rec.) & 208, 209, 223, 230 & Model Selection \\ \midrule
\textbf{DS2} & \textbf{Test} (22 rec.) & 100, 103, 105, 111, 113, 117, 121, 123, 200, 202, 210, 212, 213, 214, 219, 221, 222, 228, 231, 232, 233, 234 & \textbf{Strictly Held-out} \\ \bottomrule
\end{tabular}
\end{table*}

\subsection{Class Distribution and AAMI Mapping}
Raw annotations were mapped to the five AAMI EC57 super-classes \cite{xiao_deep_2023, luz_survey_2016} (Table~\ref{tab:class_dist}); the Q class is retained for reporting completeness but excluded from loss weighting given its extreme rarity.

\begin{table*}[!t]
\caption{Class Distribution and AAMI Standard Mapping.}
\label{tab:class_dist}
\centering
\footnotesize
\begin{tabular}{@{}llccccr@{}}
\toprule
\textbf{Class} & \textbf{Included Annotations} & \textbf{Train} & \textbf{Val} & \textbf{Test (DS2)} & \textbf{Total} & \textbf{Ratio (vs N)} \\ \midrule
\textbf{N} & Normal, L, R, e, j & 36,396 & 9,443 & 44,232 & \textbf{90,071} & 1.0 : 1 \\
\textbf{S} & A, a, J, S & 773 & 170 & 1,837 & \textbf{2,780} & 32.4 : 1 \\
\textbf{V} & V, E & 3,150 & 638 & 3,220 & \textbf{7,008} & 12.8 : 1 \\
\textbf{F} & F & 399 & 15 & 388 & \textbf{802} & 112.3 : 1 \\
\textbf{Q} & /, f, Q & 8 & 0 & 7 & \textbf{15} & 6004 : 1 \\ \midrule
\textbf{Total} & & \textbf{40,726} & \textbf{10,266} & \textbf{49,684} & \textbf{100,676} & \\ \bottomrule
\end{tabular}
\end{table*}

%% EXPERIMENTAL SETUP
\section{Experimental Setup}

\subsection{Training Details}
AdamW, initial learning rate $1\times10^{-3}$, weight decay $1\times10^{-4}$, cosine-annealed learning-rate schedule, up to 100 epochs, early stopping on validation Macro-F1 with patience 15, batch size 64. All checkpoint selection uses DS1-validation only.

\subsection{Baselines}
Table~\ref{tab:baselines} lists the four standard architectures compared against WavKAN-v2 in \S\ref{sec:results}, trained under an identical protocol, split, and epoch budget (natural-distribution sampling, class-weighted loss --- weighted cross-entropy for three of the four, and gamma-only Focal Loss for the fourth, matching each method's own standard formulation rather than stacking two imbalance-correction mechanisms, which we found in preliminary runs to cause majority-class collapse when combined).

\begin{table*}[!t]
\caption{Baseline Architectures (identical protocol/split/epoch budget as WavKAN-v2).}
\label{tab:baselines}
\centering
\footnotesize
\begin{tabular}{@{}l l l l@{}}
\toprule
\textbf{Model} & \textbf{Params} & \textbf{Description} & \textbf{Loss} \\ \midrule
ResNet1D & 62,869 & 4-block 1D residual CNN, 32 channels & Weighted CE \\
Transformer & 29,653 & $d_{model}{=}32$, 2 layers, 4 heads & Weighted CE \\
CNN+Focal & 71,061 & 4-layer 1D CNN & Focal ($\gamma{=}2$, no static weight) \\
B-Spline KAN & 118,229 & Same family as WavKAN-v2, B-spline basis, no PCWI/PWAM/RR-attention & Weighted CE \\
\textbf{WavKAN-v2} & \textbf{154,325} & \textbf{Mexican Hat, PCWI+PWAM+RR-attention, PCA curriculum} & \textbf{Focal + PCA} \\ \bottomrule
\end{tabular}
\end{table*}

%% RESULTS
\section{Results}
\label{sec:results}

\subsection{Curriculum vs. Fair Baseline (Within-Model, 20 Seeds)}
\label{sec:results_curriculum}
Table~\ref{tab:curriculum_stats} reports the paired comparison between the PCA curriculum arm and the architecturally identical fair-baseline arm, across all five AAMI classes, with Holm-Bonferroni correction applied across the five metrics tested (correcting the family-wise error rate, since testing five metrics from one comparison without correction inflates the apparent significance rate).

\begin{table*}[!t]
\caption{Curriculum vs. Fair Baseline, WavKAN-v2, 20 Seeds, DS2 Test Set. Only S-Recall survives Holm-Bonferroni correction.}
\label{tab:curriculum_stats}
\centering
\footnotesize
\begin{tabular}{@{}l c c c c c@{}}
\toprule
\textbf{Metric} & \textbf{Baseline} & \textbf{Curriculum} & \textbf{Holm-adj. $p$} & \textbf{Cohen's $d$} & \textbf{Significant?} \\ \midrule
Macro-F1 & $0.319\pm0.013$ & $0.323\pm0.018$ & $0.756$ & $-0.16$ (negl.) & No \\
V-Recall & $0.879\pm0.019$ & $0.889\pm0.021$ & $0.461$ & $-0.41$ (small) & No \\
\textbf{S-Recall} & $\mathbf{0.079\pm0.016}$ & $\mathbf{0.123\pm0.052}$ & $\mathbf{0.0016}$ & $\mathbf{-0.77}$ \textbf{(medium)} & \textbf{Yes} \\
N-Recall & $0.882\pm0.024$ & $0.865\pm0.031$ & $0.086$ & $+0.53$ (medium) & No \\
F-Recall & $0.0018\pm0.0025$ & $0.0034\pm0.0035$ & $0.461$ & $-0.32$ (small) & No \\ \bottomrule
\end{tabular}
\end{table*}

\noindent The only conclusion that survives correction is that curriculum learning provides a real, medium-effect-size improvement in S-Recall, with no statistically robust cost to any other metric. S-Recall remains low in absolute terms (0.123, i.e., roughly 88\% of S beats still missed) --- the finding is a genuine, targeted improvement, not a solution to S-class detection.

\subsection{Comparison Against Standard Architectures (Cross-Model, 20 Seeds)}
\label{sec:results_baselines}
Table~\ref{tab:baseline_comparison} and Fig.~\ref{fig:seed_stability} report WavKAN-v2 (curriculum arm) against the four baselines of Table~\ref{tab:baselines}, all under the identical protocol.

\begin{table}[!t]
\centering
\caption{Cross-Architecture Comparison, 20 Seeds Each, DS2 Test Set. All four comparisons survive Holm-Bonferroni correction.}
\label{tab:baseline_comparison}
\resizebox{\columnwidth}{!}{%
\begin{tabular}{l c c c c}
\toprule
\textbf{Model} & \textbf{Params} & \textbf{Macro-F1} & \textbf{Cohen's $d$} & \textbf{Holm $p$} \\ \midrule
ResNet1D & 62,869 & $0.351\pm0.018$ & $-1.21$ (large) & $0.0002$ \\
Transformer & 29,653 & $0.353\pm0.039$ & $-0.69$ (medium) & $0.0042$ \\
CNN+Focal & 71,061 & $0.362\pm0.015$ & $-1.39$ (large) & $0.0002$ \\
B-Spline KAN & 118,229 & $0.371\pm0.023$ & $-1.68$ (large) & $0.00004$ \\
\textbf{WavKAN-v2} & \textbf{154,325} & $\mathbf{0.323\pm0.018}$ & --- & --- \\ \bottomrule
\end{tabular}
}
\end{table}

\noindent \textbf{We report this plainly rather than reframe it}: all four baselines significantly outperform WavKAN-v2 on Macro-F1, despite WavKAN-v2 having the most parameters of the five. V-Recall, separately, is close across all five models (0.88--0.91), so the gap is specific to Macro-F1 (driven by S/F-class performance) rather than a uniform failure --- WavKAN-v2 does not lose at every axis of comparison. \S\ref{sec:discussion} discusses why, and what WavKAN-v2 offers instead of a raw Macro-F1 advantage.

\begin{figure*}[!t]
\centering
\includegraphics[width=0.85\textwidth]{baseline_comparison_seed_stability.pdf}
\caption{Cross-architecture Macro-F1 stability, 20 real seeds per model, DS2 test set. Significance markers denote paired Wilcoxon signed-rank tests against WavKAN-v2 with Holm-Bonferroni correction; effect sizes are Cohen's $d$.}
\label{fig:seed_stability}
\end{figure*}

\begin{figure*}[!t]
\centering
\includegraphics[width=\textwidth]{training_convergence_curves.pdf}
\caption{Training convergence, WavKAN-v2, fair-baseline vs.\ curriculum arm (real training-history logs, DS1 validation Macro-F1 per epoch).}
\label{fig:convergence}
\end{figure*}

\subsection{Internal Architectural Ablation}
\label{sec:results_ablation}
Table~\ref{tab:family_ablation} reports an ablation over wavelet basis and the PCWI/PWAM/RR-attention components, on both DS2 test and DS1 validation (20 seeds each; validation numbers were extracted directly from each seed's logged training history using the same selection rule used to save its checkpoint, so no additional training or test-set access was needed to obtain them).

\begin{table*}[!t]
\centering
\caption{Internal Ablation, WavKAN-v2 Family, 20 Seeds Each. Removing the RR self-attention branch yields the best Macro-F1 in the family, on both splits.}
\label{tab:family_ablation}
\footnotesize
\begin{tabular}{l c c c c c c}
\toprule
& \multicolumn{3}{c}{\textbf{DS2 Test}} & \multicolumn{3}{c}{\textbf{DS1 Validation}} \\
\textbf{Configuration} & \textbf{Macro-F1} & \textbf{V-Rec} & \textbf{S-Rec} & \textbf{Macro-F1} & \textbf{V-Rec} & \textbf{S-Rec} \\ \midrule
\textbf{Full (mexican\_hat, curriculum)} & $0.323$ & $0.889$ & $0.123$ & $0.415\pm0.020$ & $0.879\pm0.039$ & $0.305\pm0.046$ \\
No curriculum (fair baseline) & $0.319$ & $0.879$ & $0.079$ & $0.405\pm0.017$ & $0.864\pm0.041$ & $0.218\pm0.047$ \\
No PCWI & $0.329$ & $0.900$ & $0.136$ & $0.393\pm0.013$ & $0.868\pm0.040$ & $0.244\pm0.043$ \\
No PWAM & $0.310$ & $0.882$ & $0.113$ & $0.398\pm0.018$ & $0.876\pm0.031$ & $0.312\pm0.060$ \\
\textbf{No RR self-attention} & $\mathbf{0.357}$ & $0.898$ & $0.198$ & $\mathbf{0.440\pm0.014}$ & $0.884\pm0.028$ & $0.352\pm0.065$ \\
Wavelet: Morlet & $0.333$ & $0.917$ & $0.131$ & $0.422\pm0.018$ & $0.874\pm0.039$ & $0.404\pm0.066$ \\
Wavelet: DOG & $0.334$ & $0.907$ & $0.163$ & $0.404\pm0.012$ & $0.885\pm0.036$ & $0.250\pm0.051$ \\
Wavelet: B-Spline (within full config) & $0.341$ & $0.889$ & $0.273$ & $0.412\pm0.015$ & $0.863\pm0.047$ & $0.326\pm0.061$ \\ \bottomrule
\end{tabular}
\end{table*}

\noindent None of the within-family ablations reach the four external baselines' Macro-F1 range (0.351--0.371, Table~\ref{tab:baseline_comparison}). However, \textbf{removing the RR self-attention branch improves Macro-F1 from 0.323 to 0.357 and S-Recall from 0.123 to 0.198} --- the RR self-attention module, one of the architecture's headline components, is measurably net-harmful in its current form on this task, a finding we take at face value rather than omit (\S\ref{sec:discussion}). The validation-set columns are reported in the same table specifically so that the validation-test gap is visible directly: it is $\approx$0.09 for both the curriculum and no-curriculum arms, i.e., \emph{curriculum learning does not measurably reduce the generalization gap relative to the fair baseline} on this data --- we note this explicitly because an earlier internal draft of this analysis had claimed the opposite (a near-zero gap specific to curriculum learning), which we found, on directly re-deriving the validation numbers, not to be supported.

\subsection{RR-Interval Position Ablation}
\label{sec:results_rr}
Table~\ref{tab:rr_ablation} reports a leave-one-out ablation over the 5-element causal RR history (20 seeds), masking one position at a time and measuring the S-Recall delta, with a two-sided significance test per position.

\begin{table}[!t]
\centering
\caption{RR-Interval Leave-One-Out Ablation, 20 Seeds, DS2 Test Set. Baseline S-Recall (RR history intact): $0.1225\pm0.0524$.}
\label{tab:rr_ablation}
\begin{tabular}{l c c c}
\toprule
\textbf{Masked Position} & \textbf{$\Delta$S-Recall} & \textbf{$p$-value} & \textbf{Sig.?} \\ \midrule
$t-4$ & $-0.0020\pm0.0077$ & $0.081$ & No \\
$t-3$ & $-0.0016\pm0.0075$ & $0.074$ & No \\
\textbf{$t-2$} & $\mathbf{-0.0323\pm0.0492}$ & $\mathbf{6\times10^{-5}}$ & \textbf{Yes} \\
$t-1$ & $+0.0022\pm0.0095$ & $0.358$ & No \\
$t$ (current) & $-0.0045\pm0.0103$ & $0.0035$ & Yes (small) \\ \bottomrule
\end{tabular}
\end{table}

\noindent Masking the $t-2$ interval (two beats prior) causes by far the largest, most significant S-Recall degradation; the concurrent interval $t$ also reaches significance but with a much smaller effect. The raw V-Recall delta at $t-2$ is also large ($-0.0745\pm0.1132$), though we did not compute a separate significance test for V-Recall in this ablation and report it as a raw effect only. This is physiologically consistent with the compensatory-pause dynamics of ectopic beats, where the interval two beats prior carries diagnostically relevant timing information (Fig.~\ref{fig:rr_ablation}).

\begin{figure*}[!t]
\centering
\includegraphics[width=0.9\textwidth]{rr_interval_ablation.pdf}
\caption{Leave-one-out RR-interval ablation, 20 real seeds, DS2 test set. The $t-2$ position causes the largest, most significant S-Recall drop.}
\label{fig:rr_ablation}
\end{figure*}

\subsection{Cross-Dataset Zero-Shot Generalization (PTB-XL)}
\label{sec:results_ptbxl}
Table~\ref{tab:ptbxl_zero_shot} reports zero-shot evaluation (no fine-tuning) of the MIT-BIH-trained WavKAN-v2 (curriculum, seed 42) on 75,170 real PTB-XL beats. \textbf{This is a single-seed result}, not a 20-seed aggregate, and is reported as a limitation rather than a strength (\S\ref{sec:discussion}); INCART and SVDB zero-shot evaluation for WavKAN-v2 had not completed at the time of writing and are left as explicit future work rather than reported with placeholder or estimated numbers.

\begin{table}[!t]
\centering
\caption{PTB-XL Zero-Shot Generalization (Single Seed, No Fine-Tuning). Macro-F1: 0.204.}
\label{tab:ptbxl_zero_shot}
\begin{tabular}{l c c r}
\toprule
\textbf{Class} & \textbf{Recall} & \textbf{F1} & \textbf{Support} \\ \midrule
N & $0.715$ & $0.719$ & $53{,}497$ \\
S & $0.083$ & $0.098$ & $6{,}741$ \\
V & $0.176$ & $0.140$ & $5{,}889$ \\
F & --- & --- & $0$ \\
Q & $0.044$ & $0.061$ & $9{,}043$ \\ \bottomrule
\end{tabular}
\end{table}

\noindent Cross-dataset transfer is weak: Macro-F1 drops from 0.323 (DS2, in-distribution) to 0.204 (PTB-XL, zero-shot), and V-Recall in particular drops sharply (0.889 $\to$ 0.176). No Fusion-labeled beats were present in this PTB-XL extraction. We report this without softening --- it is a genuine, unresolved limitation of the current model.

\subsection{Quantified Interpretability (WAS)}
\label{sec:results_was}
To move beyond a qualitative interpretability claim, we introduce the Wavelet--ECG Alignment Score (WAS): for each learned wavelet, we compare its translation ($\mu$) and dilation ($\gamma$) parameters against literature-derived priors for QRS, P-wave, and T-wave morphology, and report a normalized alignment score per fiducial and a macro-average. WavKAN-v2 achieves $\text{WAS}_{\text{macro}} = 0.9713$ (QRS: 0.9865, P-wave: 0.9629, T-wave: 0.9644; Fig.~\ref{fig:wavelets}) --- a real, checkpoint-derived measurement, not an estimate. This metric is, by construction, unavailable to any of the four comparator architectures in Table~\ref{tab:baseline_comparison}, none of which expose an inspectable basis function.

\begin{figure*}[!t]
\centering
\includegraphics[width=\textwidth]{final_learned_wavelets.pdf}
\caption{\textbf{Learned Wavelet Bases.} WavKAN-v2 learns adaptive Mexican Hat wavelets (blue curves) whose translation/dilation parameters align with QRS/P/T-wave priors (WAS $=0.971$).}
\label{fig:wavelets}
\end{figure*}

\subsection{Deployment Efficiency and Quantization}
\label{sec:results_deploy}
Table~\ref{tab:deployment} reports real, measured CPU inference behavior (Intel x86, single-thread, no GPU) for the seed-42 curriculum checkpoint.

\begin{table}[!t]
\centering
\caption{Real Measured Deployment Metrics (CPU, FP32 vs.\ INT8 Dynamic Quantization).}
\label{tab:deployment}
\begin{tabular}{l c c}
\toprule
\textbf{Metric} & \textbf{FP32} & \textbf{INT8 (dynamic)} \\ \midrule
Model size & $0.603$ MB & $0.558$ MB ($1.08\times$) \\
Latency (batch=1, mean) & $3.18\pm1.22$ ms & $4.43\pm3.51$ ms \\
Latency (batch=64, mean) & $2.80$ ms & --- \\
Speedup vs.\ FP32 & $1.00\times$ & $0.72\times$ (\emph{slower}) \\ \bottomrule
\end{tabular}
\end{table}

\noindent We report an honest negative result: post-training dynamic INT8 quantization provides essentially no benefit for this architecture, and makes inference \emph{slower}. This is because \texttt{torch.quantize\_dynamic} only quantizes \texttt{nn.Linear} layers, while this model's parameter mass sits in the custom KAN wavelet layer and the bidirectional GRU, neither of which that method touches. A ``lightweight'' claim for this architecture should therefore rest on its real parameter count and measured FP32 latency (both genuinely favorable relative to MAK-Net's 24 ms/beat on a GPU \cite{zhao_mak-net_2025}), not on an unverified quantization benefit.

%% DISCUSSION
\section{Discussion}
\label{sec:discussion}

\subsection{An Interpretability--Accuracy Trade-off, Not an Accuracy Win}
The central honest finding of this paper is that WavKAN-v2 does not win the Macro-F1 comparison it might be expected to win (\S\ref{sec:results_baselines}): all four standard, smaller baselines significantly outperform it. We do not attribute this to an unfair comparison --- all five models share the same data, split, optimizer, and epoch budget, and each baseline uses its own standard, non-stacked imbalance-correction mechanism. We attribute it, tentatively, to the combination of (i) a wavelet-KAN parameterization that is harder to optimize at this data scale than a plain convolutional or attention stack, and (ii) the RR self-attention branch specifically adding noise rather than signal (\S\ref{sec:discussion_rr}). What the architecture offers instead is the quantified interpretability result of \S\ref{sec:results_was} (a WAS of 0.971, unavailable to any comparator) and the targeted S-Recall benefit of \S\ref{sec:results_curriculum}. We consider this trade-off --- a measured, disclosed cost in aggregate accuracy for a measured, quantified gain in structural transparency and a measured, statistically robust minority-class benefit --- to be a legitimate and clinically relevant contribution in its own right, rather than a shortfall to be explained away.

\subsection{Why the RR Self-Attention Branch Underperforms}
\label{sec:discussion_rr}
Table~\ref{tab:family_ablation} shows that removing the RR self-attention branch entirely \emph{improves} both Macro-F1 (0.323$\to$0.357) and S-Recall (0.123$\to$0.198) within the WavKAN-v2 family --- the branch is, in its current form, net-harmful. This is not inconsistent with \S\ref{sec:results_rr}'s finding that the $t-2$ RR position carries real, significant diagnostic signal for S-detection: the information is present in the raw RR history, but self-attention over a 5-element vector may dilute a sharply localized positional signal (the $t-2$ lag specifically) into an order-invariant pairwise representation, effectively discarding the one inductive bias --- ``this specific lag matters most'' --- that the raw ablation shows to be real. Concretely, WavKAN-v2's RR branch supports a \texttt{use\_rr\_attn} flag that switches between the self-attention fusion and a plain MLP over the same 5-element vector; Table~\ref{tab:family_ablation}'s ``No RR self-attention'' row \emph{is} this plain-MLP variant, already trained and evaluated across the same 20 seeds, not an untested proposal. Given that it outperforms the self-attention configuration on every metric reported, we recommend the plain-MLP RR encoding as the new default for this architecture going forward, and report the self-attention variant here as the configuration that was originally proposed and is now shown, by direct ablation, not to be the better choice.

\subsection{Augmentation Strategies for S-Class Detection}
We separately evaluated four augmentation strategies (Gaussian noise, baseline wander, a Combined strategy, and SMOTE temporal interpolation) against the DS2 test set (Fig.~\ref{fig:augmentation}; the underlying per-seed JSON for this study did not survive in the working checkout used to prepare this manuscript, so the values below are read directly from the referenced figure rather than an independently re-derivable table, and should be treated as approximate). SMOTE achieves the highest raw S-Recall ($\approx0.46$) but reduces V-Recall ($\approx0.84\to0.81$); the Combined noise strategy preserves V-Recall while modestly improving S-Recall ($\approx0.29\to0.33$). Given that V-Recall is this task's primary clinical safety endpoint, we consider the SMOTE trade-off clinically unacceptable and the Combined noise strategy the more defensible choice, consistent with Bae \& Shin's \cite{bae_handling_2025} finding that naive interpolation-based augmentation can degrade performance on imbalanced medical data.

\begin{figure}[!t]
\centering
\includegraphics[width=\columnwidth]{fig12_augmentation.pdf}
\caption{Augmentation strategy comparison, DS2 test set. SMOTE raises S-Recall but reduces V-Recall; the Combined (Gaussian + baseline-wander) strategy preserves V-Recall while modestly improving S-Recall.}
\label{fig:augmentation}
\end{figure}

\subsection{E3C Compliance}
Following the four-criterion E3C evaluation framework of Silva et al. \cite{silva_systematic_2025} (strict inter-patient evaluation; class-specific metrics under imbalance; AAMI protocol adherence; consideration of computational complexity), WavKAN-v2 satisfies all four criteria under our own re-verified compliance check, placing it among the reported 4.1\% of reviewed ECG classification studies meeting these standards \cite{silva_systematic_2025}. We note this is a protocol-adherence claim, not an accuracy claim --- it is fully compatible with \S\ref{sec:discussion}.1's disclosed Macro-F1 shortfall.

\subsection{Limitations}
\label{sec:limitations}
(1) Macro-F1 (0.323) is not only modest in absolute terms but significantly \emph{below} four standard, smaller baselines under identical conditions (\S\ref{sec:results_baselines}) --- this is the paper's most important limitation and is not softened elsewhere in this text. (2) F-class recall remains near chance ($0.0034\pm0.0035$), reflecting extreme sample scarcity (388 test beats). (3) The RR self-attention branch is measurably net-harmful in its current form (\S\ref{sec:discussion_rr}); the reported architecture therefore does not yet represent the best configuration within its own design family (Table~\ref{tab:family_ablation} shows a same-family variant without this branch scoring higher on every metric reported). (4) Zero-shot cross-dataset generalization to PTB-XL is weak and evaluated on a single seed only (\S\ref{sec:results_ptbxl}); INCART/SVDB zero-shot evaluation for this architecture was not complete at the time of writing. (5) Post-training INT8 dynamic quantization provides no benefit for this architecture (\S\ref{sec:results_deploy}); real on-device (e.g., ARM Cortex-M) latency benchmarking, as opposed to the x86 CPU measurement reported here, has not been performed. (6) The augmentation-strategy comparison (\S\ref{sec:discussion_rr}) relies on a figure whose underlying per-seed numerical data did not survive in the working checkout, so the specific values reported there are approximate.

\subsection{Future Work}
Future work will pursue: (i) adopting the already-evaluated plain-MLP RR encoding (\S\ref{sec:discussion_rr}, Table~\ref{tab:family_ablation}) as this architecture's reported default in place of the self-attention variant, since it is already shown superior on every metric and requires no new training to adopt, only a revised headline configuration; (ii) completing INCART/SVDB zero-shot evaluation and a multi-seed PTB-XL re-run; (iii) real on-device latency benchmarking on ARM-class hardware; and (iv) re-deriving the augmentation-strategy comparison from a fresh, fully-logged run so that Fig.~\ref{fig:augmentation}'s values can be reported with the same seed-level statistical rigor as the rest of this paper.

%% CONCLUSION
\section{Conclusion}
\label{sec:conclusion}
This paper introduced WavKAN-v2, a structurally interpretable, 154,325-parameter wavelet-KAN architecture for inter-patient ECG arrhythmia classification, evaluated with 20-seed, Holm-Bonferroni-corrected statistics against four standard baselines and its own architectural variants. We report two honest, load-bearing findings rather than a single headline accuracy claim: a real, statistically robust, targeted improvement in Supraventricular recall from Progressive Curriculum Anchoring (0.079$\to$0.123, Holm $p=0.0016$), and a real, statistically robust Macro-F1 \emph{disadvantage} relative to four smaller standard architectures. In place of an accuracy advantage, WavKAN-v2 offers a quantified interpretability metric (WAS $=0.971$) unavailable to any of its comparators, real measured deployment latency, and an honestly reported negative quantization result. We further identify, via our own ablation, that the architecture's RR self-attention branch is currently net-harmful and should be revisited before this architecture is treated as finished. We believe this level of disclosed, statistically corrected reporting --- including the results that do not favor the proposed model --- is itself a needed contribution to a literature in which inter-patient comparisons are frequently reported without correction for multiple comparisons or a matched training recipe across baselines.

\section*{Acknowledgment}
The authors acknowledge the Vellore Institute of Technology (VIT) for funding the Article Processing Charge (APC) of this manuscript. The authors declare no competing financial interests.

\textbf{CRediT Author Statement:}
\textbf{Venkateramanan Manivannan:} Conceptualization, Methodology, Software, Validation, Formal Analysis, Investigation, Data Curation, Writing -- Original Draft, Visualization.
\textbf{Ramanathan Lakshmanan:} Conceptualization, Resources, Validation, Supervision, Project Administration, Funding Acquisition, Writing -- Review \& Editing.

\textbf{Data Availability:} The datasets analyzed during the current study (MIT-BIH Arrhythmia Database) are publicly available in the PhysioNet repository: \url{https://physionet.org/content/mitdb/1.0.0/} \cite{moody_mit-bih_1992}.

\textbf{Code Availability:} Source code and pre-trained weights are available at: \url{https://github.com/vrhsr/WavKAN-CL}.
\bibliographystyle{IEEEtran}
\bibliography{references}

\begin{IEEEbiography}[{\includegraphics[width=1in,height=1.25in,clip,keepaspectratio]{venkateramanan.jpg}}]{Venkateramanan Manivannan}
received the B.E. degree in Computer Science and Engineering and is currently pursuing the M.Tech degree in Computer Science and Engineering at the School of Computer Science and Engineering (SCOPE), Vellore Institute of Technology (VIT), Vellore, India.

His research interests include machine learning, deep learning, artificial intelligence, neural network architectures, explainable AI, natural language processing, computer vision, biomedical signal processing, and computationally efficient models for edge deployment. His current work focuses on developing interpretable and resource-efficient deep learning frameworks for healthcare and safety-critical applications.
\end{IEEEbiography}

\begin{IEEEbiography}[{\includegraphics[width=1in,height=1.25in,clip,keepaspectratio]{ramanathan.jpg}}]{Ramanathan Lakshmanan}
received the B.E. degree in Computer Science and Engineering from Bharathidasan University, Tiruchirappalli, India, the M.E. degree in Computer Science from Sathyabama University, Chennai, India, and the Ph.D. degree in Computer Science and Engineering from VIT University, Vellore, India. He is currently a Professor with the Department of IoT, School of Computer Science and Engineering, VIT, Vellore, India, where he has over 18 years of teaching and research experience.

His research interests include data science, cyber-physical systems, fog and edge computing, cloud computing, big data classification, IoT, artificial intelligence, and machine learning. His ongoing research focuses on predictive analytics, classification, clustering, and optimization. He serves as an Associate Editor for the Journal of Network: Computation in Neural Systems, and is an Editorial Advisory Board Member for the Journal of Integrated Science and Technology. He is an active member of IACSIT, CSI, ACM, IEEE (WIE), and ACEEE.
\end{IEEEbiography}

\end{document}
